\section{Related Work}\label{sec:related}

Prior research has been conducted into how LLMs sort their content compared to media outlets. The question of whether LLMs deprioritize content or merely synthesize is explored by Li and Taneja~\cite{li2026mirror}, who conducted an audit and found that LLMs such as ChatGPT and Gemini deprioritize contentious topics compared to traditional news outlets. Similarly, Minici et al.~\cite{minici2025auditing} audited LLM editorial biases, using algorithmic auditing to examine LLMs taking on the role of news outlets and finding that LLMs use fewer unique outlets than Google News. We are similarly interested in the scope of news outlets that LLMs pull from and how LLMs may curate information. Our work can add to this research by filling in the gap on how paid/tiered models contribute to the deprioritization of content. We will also investigate a more user-oriented perspective of traditional outlets and LLM outputs, using news outlets as a ``baseline'' for what we may consider credible news.

Additional research has been conducted into how LLMs convey news information. For example, Khatiwada et al.~\cite{khatiwada2026rewrites} examine how variations in tone and framing affect readers' judgments of credibility and their perception of information. Their findings showed that extreme sentiment lowered trustworthiness. While they explored the effects of LLM-modified articles, we will go into further depth on the differences within unmodified local news across different outlets.

Another area of research is LLM information sourcing. For instance, Zhang et al.~\cite{zhang2025institutional} examine LLMs within journalistic workflows and how incorporating LLMs into journalism can speed up production at the cost of citing fewer sources. This area is more focused on the role of journalists and how they may incorporate LLMs into newsroom practices. It aligns with our study less, but it can be useful for theorizing future integrations in the opposite direction, in other words, if LLMs adopt customs of ``local news'' to cater to audiences.